% Experimental setup.
\section{Experimental Setup}
\label{sec:setup}

\subsection{Data and cross-condition protocols}
\label{sec:data}

We use the CWRU bearing data \cite{zheng2019cwru}, which provides drive-end
acceleration under four motor loads (0--3\,hp) and ten classes (normal, inner-race,
ball, and outer-race faults at three severities each); load is the cross-domain
axis, shifting rotation speed and fault impulsiveness. Two disjoint protocols
carry the main evaluation:
\begin{itemize}
  \item \textbf{cwru\_03}: single-source hardest case. Train on CWRU load 0,
  transfer to load 3, ten classes.
  \item \textbf{cwru\_0123}: multi-source case. Train on loads 0,1,2, transfer
  to load 3, ten classes.
\end{itemize}
Both are recording- and condition-disjoint (Sec.~\ref{sec:method}), so no test
window shares a recording or operating condition with training.

\subsection{Models and implementation}
\label{sec:models}

We compare the raw-signal \textbf{1D-CNN} baseline \cite{ince2016real,
wen2018new,zhang2017new}, the source-only structural baseline
\textbf{DynHGNN+}, the full \textbf{GA-DCHG}, and its ablations, namely
GA-DCHG $-$align, $-$temporal, and $-$anchor. Embedding dimension $m=3$ chosen
by lag through mutual information \cite{fraser1986}, segments of $S=64$ state
points,
$n_{\mathrm{seg}}=4$ per window, neighbourhood order $k=8$, maximal hyperedge
size $K_{\max}=12$, anchor strength $\eta=1.0$, propagation $\tau=0.5$, and
alignment weight $\lambda=0.5$; we train for $60$ epochs with batch size $16$,
Adam with learning rate $2\times10^{-3}$, and the alignment term active from the
first epoch. Checkpoints are selected by source validation accuracy; every
protocol runs over three seeds $\{0,1,2\}$ and we report mean $\pm$ std
target-test accuracy.
